\clearpage\begingroup\SourceLeaf{83a}\HistoricalHead{78}{L}{PHYSIK. BRIEFWECHSEL.}
\begin{center}[5.]\end{center}
\begin{center}\textsc{Gauss} an \textsc{Encke}.\quad Göttingen, 12. Mai 1832.\end{center}
\noindent . . . . . . . . . . . . . . . . . . . . .\par
\indent Seit Anfang dieses Jahres habe ich mich sehr viel mit dem Erdmagnetismus
beschäftigt, sowohl von theoretischer als praktischer Seite. In letzterer
Beziehung ist es mir schon jetzt gelungen, durch Anwendung neuer Einrichtungen,
die sich durch Einfachheit, Sicherheit, Schärfe — und Wohlfeilheit
empfehlen — meine Erwartungen nicht bloss erfüllt, sondern weit übertroffen
zu sehen; die Messungen, die sich auf die absolute Declination, deren Intensität,
und die täglichen und stündlichen Variationen von beiden beziehen,
erhalten eine Genauigkeit, die der der astronomischen Beobachtungen fast
gleich kommt. Auch die Zurückführung der Intensität auf eine absolute Einheit
denke ich zu einer vergleichungsweise grossen Schärfe zu bringen und
späterhin mich auch mit der Inclination zu beschäftigen, wozu aber sehr gut
gearbeitete Aufhängungsaxen (\textls[120]{aber sonst nichts}) erforderlich sein werden.
Bald hoffe ich diese Dinge zu einer solchen Reife zu bringen, dass ich öffentlich
etwas darüber bekannt machen kann. Einstweilen habe ich am 4. u. 5.
Mai viele Aufzeichnungen der täglichen Variation gemacht; in Zukunft soll
diess aber vollständiger geschehen.

Aus den Zeitungen sehe ich, dass H[err] von \textsc{Humboldt} aus Paris zurückgerufen
sei; ist dies gegründet und ist er schon in Berlin angekommen? Ich
würde dann mit Vergnügen ihm einige vorläufige Mittheilungen von jenen
Beschäftigungen machen. . . . . . . . . . . . . . . . . . . . . .
\clearpage\endgroup
